\scanpage{39}
\margin{\handwritten{SGA $\overline{\underline{4}}$ $\overline{\underline{\mathrm{VIII}}}$}}

10. \uline{Compléments}: \uline{espaces algébriques et étendues algébriques}.

Nous donnons sous forme d'exercices quelques compléments, qui ne seront pas utilisés dans la suite du Séminaire.

10.1. Soit $(X,\underline{O}_X)$ un schéma, \deleted{notons par posons $X_{\mathrm{ét}}=\operatorname{Top}(X)$, de} sorte que\deleted{\unsure{n|r}} $(X_{\mathrm{ét}}^{\handwritten{\sim}},\underline{O}_X)$ est un topos annelé. Montrer qu'on peut réo\unsure{\deleted{n}o|\deleted{o}n}ntsr\unsure{\deleted{y}u|\deleted{u}y}ire le schéma $X$ à isomorphique unique près, quand on connait le topos annelé $(X_{\mathrm{et}}^{\handwritten{\sim}},\underline{O}_X)$ à équivalence près. (Utiliser 6.1 pour rédonstituer \deleted{$\operatorname{Top}(X)$} le topos $\operatorname{Top}(X)=\operatorname{Ouv}(X_{\mathrm{ét}}^{\handwritten{\sim}})^{\handwritten{\sim}}$, puis utiliser IV 4.\handwritten{2}.\handwritten{3}.) Plus \figure{\begin{tikzpicture}[x=1cm,y=1cm]\draw (0,0.70) -- (0.27,0.70) -- (0.29,0.75);\draw (0.12,0) -- (0.14,0.22);\end{tikzpicture}} précisément, montrer que si $X,X'$ sont deux schémas, alors l'applic\deleted{\unsure{a|c}}\unsure{\deleted{t}a|\deleted{a}t}tion naturelle
\[
\operatorname{Hom}(X,X')\handwritten{\longrightarrow}\operatorname{Homtoplocan}(X_{\mathrm{ét}}^{\handwritten{\sim}},X'_{\mathrm{ét}}{}^{\handwritten{\sim}})
\]
\margin{\handwritten{\unsure{$\chi$|$\times$}}}
(notations de IV \handwritten{11.}) induit une équivalence entre la catégorie discrète définie par l'ensemble $\operatorname{Hom}(X,X')$ et la catégorie \uline{Homtoplocan}$(X_{\mathrm{ét}}^{\handwritten{\sim}},X'_{\mathrm{ét}}{}^{\handwritten{\sim}})$ (où pour simplifier on désigne par $X_{\mathrm{ét}}^{\handwritten{\sim}}$ le topos étale de $X$ considéré comme topos annelé). (Utiliser IV 4.\handwritten{2}.\handwritten{3}.\deleted{\unsure{et}}\handwritten{)}
\margin{\handwritten{\deleted{\unsure{etc.}}}}

10.2. Soit $(E,\underline{O}_E)$ un topos annelé. Montrer que les conditions suivantes sont équivalentes (comparer IV \handwritten{9.}):
\margin{\handwritten{$\times$}}

(i) Il existe une famille $(X_i)_{i\handwritten{\in}I}$ d'objets de $E$ couvrant l'objet final, et une famille $(X'_i)_{i\handwritten{\in}I}$ de schémas, tels que pour tout $i\handwritten{\in}I$, le topos annelé induit $E_{/X_i}$ soit équivalent à $X'_{i\,\mathrm{ét}}{}^{\handwritten{\sim}}$.

(ii) Comme (i), avec $I$ réduit à un élément.

(iii) Il existe un schéma $X_{\handwritten{0}}$, et une prérelation d'équivalence $\deleted{R}\handwritten{X_1}$ dans $X_{\handwritten{0}}$ au sens de la catégorie (Sch) des schémas \handwritten{,}\deleted{(\illegible{2})}\unsure{,} telle que $R$ soit étale i.e. le morphisme $p_1:\deleted{R}\handwritten{X_1\to}X_{\handwritten{0}}$ est étale, et que le topos annelé $(E,\underline{O}_E)$ soit équivalent au topos annelé déduit par recollement du diagramme \inserted{\handwritten{$\frac12$-simplicial}} de topos annelés \inserted{\handwritten{étales}} \handwritten{associé\deleted{\unsure{s|e}} au diagramme de schémas}
\figure{\begin{tikzpicture}[x=.025cm,y=.025cm]
\draw[line width=.65pt] (0,0) .. controls (-2,27) and (4,34) .. (22,39) .. controls (68,47) and (153,49) .. (184,54) .. controls (190,56) and (191,60) .. (191,65);
\draw[line width=.65pt] (319,0) .. controls (316,25) and (325,32) .. (345,38) .. controls (380,47) and (426,49) .. (438,57) .. controls (444,62) and (446,68) .. (446,73);
\end{tikzpicture}}
\diagram{\begin{tikzcd}[ampersand replacement=\&]
X_{\handwritten{2}}\arrow[r,shift left=1.2ex]\arrow[r]\arrow[r,shift right=1.2ex] \& X_{\handwritten{1}}\arrow[r,shift left=.6ex]\arrow[r,shift right=.6ex] \& X_o
\end{tikzcd}}
% archive sheet number: 38
